% Options for packages loaded elsewhere
\PassOptionsToPackage{unicode}{hyperref}
\PassOptionsToPackage{hyphens}{url}
\documentclass[
  ignorenonframetext,
]{beamer}
\newif\ifbibliography
\usepackage{pgfpages}
\setbeamertemplate{caption}[numbered]
\setbeamertemplate{caption label separator}{: }
\setbeamercolor{caption name}{fg=normal text.fg}
\beamertemplatenavigationsymbolsempty
% remove section numbering
\setbeamertemplate{part page}{
  \centering
  \begin{beamercolorbox}[sep=16pt,center]{part title}
    \usebeamerfont{part title}\insertpart\par
  \end{beamercolorbox}
}
\setbeamertemplate{section page}{
  \centering
  \begin{beamercolorbox}[sep=12pt,center]{section title}
    \usebeamerfont{section title}\insertsection\par
  \end{beamercolorbox}
}
\setbeamertemplate{subsection page}{
  \centering
  \begin{beamercolorbox}[sep=8pt,center]{subsection title}
    \usebeamerfont{subsection title}\insertsubsection\par
  \end{beamercolorbox}
}
% Prevent slide breaks in the middle of a paragraph
\widowpenalties 1 10000
\raggedbottom
\AtBeginPart{
  \frame{\partpage}
}
\AtBeginSection{
  \ifbibliography
  \else
    \frame{\sectionpage}
  \fi
}
\AtBeginSubsection{
  \frame{\subsectionpage}
}
\usepackage{iftex}
\ifPDFTeX
  \usepackage[T1]{fontenc}
  \usepackage[utf8]{inputenc}
  \usepackage{textcomp} % provide euro and other symbols
\else % if luatex or xetex
  \usepackage{unicode-math} % this also loads fontspec
  \defaultfontfeatures{Scale=MatchLowercase}
  \defaultfontfeatures[\rmfamily]{Ligatures=TeX,Scale=1}
\fi
\usepackage{lmodern}
\ifPDFTeX\else
  % xetex/luatex font selection
\fi
% Use upquote if available, for straight quotes in verbatim environments
\IfFileExists{upquote.sty}{\usepackage{upquote}}{}
\IfFileExists{microtype.sty}{% use microtype if available
  \usepackage[]{microtype}
  \UseMicrotypeSet[protrusion]{basicmath} % disable protrusion for tt fonts
}{}
\makeatletter
\@ifundefined{KOMAClassName}{% if non-KOMA class
  \IfFileExists{parskip.sty}{%
    \usepackage{parskip}
  }{% else
    \setlength{\parindent}{0pt}
    \setlength{\parskip}{6pt plus 2pt minus 1pt}}
}{% if KOMA class
  \KOMAoptions{parskip=half}}
\makeatother
\setlength{\emergencystretch}{3em} % prevent overfull lines
\providecommand{\tightlist}{%
  \setlength{\itemsep}{0pt}\setlength{\parskip}{0pt}}
\usepackage{bookmark}
\IfFileExists{xurl.sty}{\usepackage{xurl}}{} % add URL line breaks if available
\urlstyle{same}
\hypersetup{
  hidelinks,
  pdfcreator={LaTeX via pandoc}}

\author{\texorpdfstring{}{}}
\date{}

\begin{document}

\section{Program graph and formal
foundations}\label{program-graph-and-formal-foundations}

\section{Methodology: Program Graphs and Intermediate
Representations}\label{methodology-program-graphs-and-intermediate-representations}

\begin{frame}[fragile]{Program graph}
\protect\phantomsection\label{program-graph}
Let \(G = (V, E)\) denote a directed graph whose vertices \(V\)
represent program entities and whose edges \(E\) represent
relationships. COGANT assigns each node a \textbf{kind} (for example
function, variable, type) and a \textbf{semantic role} (for example
definition versus use). Each node and edge may carry:

\begin{itemize}
\tightlist
\item
  A \textbf{stable identifier} for persistence across runs when hashing
  allows.
\item
  \textbf{Type strings} and attributes recovered from the front end.
\item
  A \textbf{confidence} score in \([0,1]\) and \textbf{provenance}
  metadata (source code, type system, control flow, heuristic, external
  tool).
\end{itemize}

This follows the same philosophical line as code property graphs that
fuse AST, control flow, and data flow into one analyzable structure
{[}@yamaguchi2014modeling{]}, but orients the representation toward
\textbf{Generalized Notation Notation (GNN) export} and, optionally,
tensorized views of the same graph: node kinds and roles map to discrete
indices in both forms, and optional embeddings can augment features as
described in \texttt{../cogant/docs/export/README.md}.

\begin{block}{Structural equivalence}
\protect\phantomsection\label{structural-equivalence}
Two program graphs \(G_1 = (V_1, E_1)\) and \(G_2 = (V_2, E_2)\) are
usefully compared via typed graph isomorphism: a bijection
\(\phi : V_1 \to V_2\) preserving edge structure and relevant labels
supports deduplication and cross-repository linking when interfaces
align.

\begin{equation}
\label{eq:typed-iso}
(u, v) \in E_1 \iff (\phi(u), \phi(v)) \in E_2
\end{equation}

Equation \ref{eq:typed-iso} formalizes the structural invariant we check
during deduplication and cross-repository linking: a candidate bijection
\(\phi\) is accepted only when every edge in \(G_1\) has a corresponding
edge between the images of its endpoints in \(G_2\) (and vice versa), so
that label-preserving matches strictly preserve the adjacency structure
of both program graphs.
\end{block}
\end{frame}

\begin{frame}[fragile]{Formal definitions}
\protect\phantomsection\label{formal-definitions}
This subsection makes the objects manipulated by the pipeline
mathematically explicit. The definitions are stated for the shipped
v0.5.x engine; where the implementation differs from the general form
(for example because an edge kind is not yet emitted by the Python front
end), the difference is noted with a forward reference to Section 6.

\textbf{Definition 1 (Program graph).} A \textbf{program graph} is a
tuple \(G = (V, E, \lambda_V, \lambda_E, \tau)\) where

\begin{itemize}
\tightlist
\item
  \(V\) is a finite set of program nodes (modules, classes, methods,
  functions, and --- on languages where the front end emits them ---
  variables and control-flow sites);
\item
  \(E \subseteq V \times V \times K\) is a finite set of typed directed
  edges drawn from the edge-kind alphabet
  \(K \supseteq \{\text{READS}, \text{WRITES}, \text{CALLS}, \text{CONTAINS}, \text{INHERITS}, \text{IMPORTS}, \text{OBSERVES}, \text{MUTATES}, \text{DEPENDS\_ON}\}\);
\item
  \(\lambda_V : V \to \mathcal{N}\) labels each node with a node kind
  \(\mathcal{N} \supseteq \{\text{MODULE}, \text{CLASS}, \text{METHOD}, \text{FUNCTION}, \text{CONFIGURATION}, \ldots\}\);
\item
  \(\lambda_E : E \to K\) is the (trivial) projection onto the edge
  kind;
\item
  \(\tau : V \to (T \cup \{\bot\})\) maps each node to a type annotation
  recovered from the front end or to \(\bot\) when no annotation is
  available.
\end{itemize}

The shipped Python front end populates the kinds
\(\{\text{MODULE}, \text{CLASS}, \text{METHOD}, \text{FUNCTION}\}\) and
the edge kinds
\(\{\text{CALLS}, \text{CONTAINS}, \text{READS}, \text{WRITES}, \text{IMPORTS}, \text{INHERITS}\}\);
the remaining kinds in \(\mathcal{N}\) and \(K\) are declared in
\texttt{cogant.schemas.core} but currently emitted only when other
parsers or dynamic enrichment stages fire. The empirical distribution on
the six packaged fixtures is recorded in Tables 4 and 5 of Section 6.

\textbf{Definition 2 (Translation rule).} A \textbf{translation rule} is
a quadruple \(r = (\varphi_r, \kappa_r, w_r, p_r)\) where

\begin{itemize}
\tightlist
\item
  \(\varphi_r : \mathcal{G} \to 2^{\mathcal{F}}\) is a computable
  predicate that, given a program graph \(G \in \mathcal{G}\), returns a
  (possibly empty) set of matched fragments \(\mathcal{F}\) --- each
  fragment a finite tuple of node ids and optional edge ids;
\item
  \(\kappa_r \in \mathcal{K}_M\) is the mapping kind the rule assigns on
  success, drawn from the mapping-kind alphabet
  \(\mathcal{K}_M = \{\text{HIDDEN\_STATE}, \text{OBSERVATION}, \text{ACTION}, \text{POLICY}, \text{CONSTRAINT}, \text{PREFERENCE}, \text{CONTEXT}, \text{DATA\_FLOW}, \text{ERROR\_HANDLING}, \text{CIRCUIT\_BREAKER}, \text{ORCHESTRATION}\}\);
\item
  \(w_r \in (0, 1]\) is the base confidence weight attached to mappings
  produced by \(r\);
\item
  \(p_r \in \mathbb{Z}\) is the rule priority consulted during conflict
  resolution (Algorithm 2).
\end{itemize}

Concretely, each \texttt{TranslationRule} subclass in
\texttt{../cogant/py/cogant/translate/rules/} exposes \(\varphi_r\) as
the \texttt{matches(graph,\ query)} method, encodes \(\kappa_r\) in the
\texttt{mapping\_kind} property, embeds \(w_r\) in the
\texttt{confidence\_score} field of the returned
\texttt{SemanticMapping}, and exposes \(p_r\) through the
\texttt{priority} property. The nineteen shipped rules span five
families: five structural rules (\texttt{ReadOnlyInput},
\texttt{MutatingSubsystem}, \texttt{Inheritance}, \texttt{Containment},
\texttt{DataPipeline}), five semantic rules (\texttt{Observation},
\texttt{Action}, \texttt{Policy}, \texttt{Preference},
\texttt{Context}), two control rules (\texttt{Config},
\texttt{FeatureFlag}), three behavioural rules (\texttt{Orchestrator},
\texttt{TestAssertion}, \texttt{EventBus}), and four resilience rules
(\texttt{RetryPattern}, \texttt{ErrorBoundary},
\texttt{SingletonAccess}, \texttt{CircuitBreaker}).

\textbf{Definition 3 (Fixpoint semantics).} Let \(\mathcal{M}\) be the
(finite) set of all possible semantic mappings that any finite
composition of rules in \(R\) could emit on a fixed graph \(G\). Define
the rule-application operator
\(F_{G,R} : 2^{\mathcal{M}} \to 2^{\mathcal{M}}\) by \begin{equation}
\label{eq:fixpoint-operator}
F_{G,R}(S) \;=\; S \,\cup\, \bigl\{\, r.\text{apply}(G, f) \,\bigm|\, r \in R,\ f \in \varphi_r(G),\ r.\text{apply}(G, f) \neq \bot \,\bigr\}.
\end{equation} The \textbf{translation} of \(G\) under \(R\) is the
least fixpoint \begin{equation}
\label{eq:least-fixpoint}
T^{*}(G) \;=\; \bigsqcup_{k \geq 0} F_{G,R}^{k}(\emptyset) \;=\; \lim_{k \to \infty} F_{G,R}^{k}(\emptyset),
\end{equation} computed iteratively by
\texttt{TranslationEngine.translate()} and then post-processed by
\texttt{\_resolve\_conflicts()} (Algorithm 2). The post-processing step
applies an anti-monotone pruning to the fixpoint, and is therefore
specified outside of \(F_{G,R}\) rather than folded into it.

\textbf{Definition 4 (Markov blanket partition).} Given a program graph
\(G = (V, E, \lambda_V, \lambda_E, \tau)\) and a \textbf{seed set}
\(S \subseteq V\) selected by one of the five strategies in
\texttt{MarkovBlanketExtractor} (\texttt{explicit}, \texttt{module},
\texttt{kind}, \texttt{auto}, \texttt{mapping\_kind}), the
\textbf{Markov blanket partition}
\(\Pi_{G, S} : V \to \{\mu, s, a, \eta\}\) is defined on the undirected
projection of \(G\) as follows. Let
\(N^{\text{in}}(v) = \{u : (u, v, k) \in E\}\) and
\(N^{\text{out}}(v) = \{u : (v, u, k) \in E\}\). Then \begin{equation}
\label{eq:markov-partition}
\Pi_{G,S}(v) \;=\;
\begin{cases}
\mu & \text{if } v \in S \text{ and } (N^{\text{in}}(v) \cup N^{\text{out}}(v)) \subseteq S, \\
a   & \text{if } v \in S \text{ and } N^{\text{out}}(v) \setminus S \neq \emptyset, \\
s   & \text{if } v \in S \text{ and } N^{\text{out}}(v) \subseteq S \text{ and } N^{\text{in}}(v) \setminus S \neq \emptyset, \\
\eta & \text{if } v \notin S.
\end{cases}
\end{equation} The four-way case split is realised verbatim by
\texttt{partition\_by\_seeds()} in
\texttt{../cogant/py/cogant/markov/blanket.py}: the function precomputes
the in/out adjacency of every node in \(O(|V| + |E|)\) time via
\texttt{\_bidirectional\_adjacency()}, then walks \(V\) once and assigns
each node to exactly one of \(\mu\), \(s\), \(a\), \(\eta\).
Bidirectional boundary nodes (those with both external in- and
out-neighbours) are assigned to \(a\) by convention, and are
additionally tagged with a \texttt{bidirectional} metadata flag so that
downstream consumers can recover the distinction.

\textbf{Definition 5 (A/B/C/D matrices).} Given a translation
\(T^{*}(G)\) and a compiled state-space
\((\mathbf{V}, \mathbf{O}, \mathbf{A})\) of hidden-state variables,
observation modalities, and actions, the \textbf{generative-model
matrices} of COGANT are

\begin{align}
\label{eq:matrices-defn}
A &\in \mathbb{R}^{|\mathbf{O}| \times |\mathbf{V}|},    &A_{ij} &= P(o_i \mid s_j), & \sum_i A_{ij} &= 1, \\
B &\in \mathbb{R}^{|\mathbf{V}| \times |\mathbf{V}| \times |\mathbf{A}|}, & B_{i j k} &= P(s'_i \mid s_j, a_k), & \sum_i B_{ijk} &= 1, \\
C &\in \mathbb{R}^{|\mathbf{O}|}, & C_i &= \log \tilde{P}(o_i), & \\
D &\in \mathbb{R}^{|\mathbf{V}|}, & D_j &= P(s_j \mid t = 0), & \sum_j D_j &= 1.
\end{align}

\(A\) is derived from
\(\{\text{READS}, \text{OBSERVES}, \text{DEPENDS\_ON}\}\) edges between
observation and hidden-state nodes; \(B\) is derived from
\(\{\text{WRITES}, \text{MUTATES}\}\) edges from action to hidden-state
nodes, with identity fallback when an action writes nothing; \(C\) is
derived from the signed confidence scores of CONSTRAINT/PREFERENCE
mappings adjacent to each observation; and \(D\) is derived from
CONFIGURATION-neighbour bias on hidden-state variables or falls back to
a uniform prior. The derivation is performed by \texttt{GNNMatrices} in
\texttt{../cogant/py/cogant/gnn/matrices.py} and uses no numerical
dependencies beyond the Python standard library. Rows (for \(A\)) and
columns (for \(B\)) and the vectors \(D\) are normalised to valid
probability distributions using the high-direct / low-indirect mass
defaults \((0.9, 0.1)\) imported from the upstream PyMDP placeholder
convention; the \(C\) vector is a log-preference and is not normalised.

\begin{block}{Theorems}
\protect\phantomsection\label{theorems}
\textbf{Theorem 1 (Fixpoint termination).} Let
\(G = (V, E, \lambda_V, \lambda_E, \tau)\) be a finite program graph
with \(|V| = n\), let \(R\) be a finite rule set with \(|R| = k\), and
let \(F_{G,R}\) be the rule-application operator of Definition 3. Then
the Kleene chain \begin{equation}
\label{eq:kleene-chain}
\emptyset \;\subseteq\; F_{G,R}(\emptyset) \;\subseteq\; F_{G,R}^{2}(\emptyset) \;\subseteq\; \cdots
\end{equation} stabilises in at most \(|\mathcal{M}|\) iterations, where
\(|\mathcal{M}| \leq n \cdot |\mathcal{K}_M|\) is an upper bound on the
number of distinct mapping ids any rule set can produce on \(G\). In
particular, the shipped engine with its default cap \(K = 10\) converges
on every fixture in \(\{\)\texttt{calculator}, \texttt{event\_pipeline},
\texttt{flask\_mini}, \texttt{flask\_app}, \texttt{requests\_lib},
\texttt{json\_stdlib}\(\}\) within the cap.

\emph{Proof sketch.} Each application step either adds at least one new
mapping id to the accumulating set \(S\) or leaves \(S\) unchanged; the
engine's explicit \texttt{mapping.id\ not\ in\ self.mappings} guard
(\texttt{engine.py} lines 299--303) is what enforces the
monotone-plus-bounded invariant. Because \(F_{G,R}\) is monotone on the
finite lattice \((2^{\mathcal{M}}, \subseteq)\) and bounded above by
\(\mathcal{M}\), the Kleene chain is an ascending chain in a finite
lattice and therefore stabilises in at most \(|\mathcal{M}|\) steps. The
stabilisation point is, by construction, the least fixpoint of
\(F_{G,R}\) above \(\emptyset\). The worst-case bound
\(n \cdot |\mathcal{K}_M|\) is attained when every node receives at most
one mapping of each kind, which is the maximum any rule set can inject
before the conflict-resolution step prunes overlaps. Empirically every
packaged fixture converges in a single pass because the shipped rules
are disjoint on the node kinds they target: Tables 4 and 6 of Section 6
record \texttt{mappings\_total} values that equal the sum of per-kind
counts, which is only possible if no node was touched twice.
\(\blacksquare\)

\textbf{Theorem 2 (Markov blanket completeness).} For any program graph
\(G = (V, E, \lambda_V, \lambda_E, \tau)\) with \(V \neq \emptyset\) and
any seed set \(S \subseteq V\), the partition \(\Pi_{G, S}\) of
Definition 4 is \textbf{total} (every \(v \in V\) receives exactly one
role) and \textbf{mutually exclusive} (no \(v\) is assigned two roles).

\emph{Proof sketch.} Totality: the case split in Equation
\ref{eq:markov-partition} exhausts the boolean product
\((v \in S) \times (N^{\text{out}}(v) \setminus S = \emptyset) \times (N^{\text{in}}(v) \setminus S = \emptyset)\).
If \(v \notin S\) the fourth case applies. If \(v \in S\) there are
three subcases depending on whether \(v\) has external out-neighbours,
external in-neighbours only, or no external neighbours at all; the
engine resolves the ``both external in and external out'' subcase to
\texttt{ACTIVE} by convention, matching lines 205--211 of
\texttt{blanket.py}. Every \(v \in V\) therefore matches exactly one
branch, so \(\Pi\) is a total function. Mutual exclusivity: the branches
are pairwise disjoint because they condition on mutually exclusive
predicates on
\((v \in S, N^{\text{in}}(v) \setminus S, N^{\text{out}}(v) \setminus S)\).
The implementation materialises the four role sets \texttt{internal},
\texttt{sensory}, \texttt{active}, \texttt{external} as disjoint Python
sets and writes into exactly one on each iteration of the main loop, so
the in-memory invariant matches the mathematical one. \(\blacksquare\)

\textbf{Theorem 3 (Matrix validity).} If \(|\mathbf{O}| \geq 1\),
\(|\mathbf{V}| \geq 1\), and \(|\mathbf{A}| \geq 1\), then the matrices
\((A, B, C, D)\) produced by \texttt{GNNMatrices.compute\_A/B/C/D}
satisfy the stochastic conditions of Definition 5 within a numerical
tolerance of \(10^{-6}\).

\emph{Proof sketch.} Rows of \(A\) are produced in Equation
(\ref{eq:matrices-defn}) with explicit normalisation
(\texttt{\_normalize\_row()} at \texttt{matrices.py} line 277), which
divides by the row sum when it exceeds \(\varepsilon = 10^{-9}\) and
returns a uniform row otherwise. Columns of each action slice
\(B[:,:,k]\) are produced by the same helper on lines 334--348, with an
identity fallback ensuring the resulting column is never all-zero. The
prior \(D\) is built as a weighted vector of confidence scores and
passed through \texttt{\_normalize\_vector()}, again with a uniform
fallback when all weights are below \(\varepsilon\). The implementation
therefore establishes the sum-to-one invariant by construction, and
\texttt{validate\_shapes()} (lines 554--603) enforces a tolerance of
\(10^{-6}\) on every computed matrix before the pipeline accepts the
bundle; all six packaged fixtures pass \texttt{GNNValidator} with zero
errors (Table 7), which is observationally the same statement as Theorem
3 on those fixtures. The \(C\) vector is a log-preference and is exempt
from the sum-to-one requirement. \(\blacksquare\)
\end{block}
\end{frame}

\end{document}
